% Post-CP20 conditional argument. No assertion that its extra hypothesis
% holds on the actual nonordinary Cartier locus.
\subsection{A conditional polar-tensor test for the remaining Lang locus}
\status{Proved under the explicit extra stability hypothesis; independently checked.}

\begin{lemma}[A cubic containing a Frobenius-defined smooth sixfold]
\label{lem:polar-third-derivative-Lang-candidate}
Let $k$ be algebraically closed of characteristic two, let $V=k^9$, and
let $f$ be a homogeneous cubic on $V$. Suppose that on an open subset
of $V$ there are three rational vector-valued functions $k_i(x^{[2]})$
which are everywhere linearly independent. Set
\[
                 g_i(x)=x^t k_i(x^{[2]}),\qquad i=1,2,3.
\]
If an irreducible six-dimensional component of $g_1=g_2=g_3=0$
is contained in $f=0$, the alternating polar three-form of $f$
is not stable for $\operatorname{SL}(V)$.
\end{lemma}
\begin{proof}
Shrink to an open on which all functions are regular. Ordinary
coordinate differentiation gives
\[
 dg_i=k_i(x^{[2]})^t dx,\qquad D_vD_wg_i=0
\]
for all constant vector fields $v,w$, because all coefficients of
$k_i(x^{[2]})$ have zero differential, including rational denominators.
The three independent differentials make their zero scheme smooth
of dimension six. At a general point $z$ of the indicated component,
its local ideal is therefore generated by $g_1,g_2,g_3$; after shrinking
again one can write $f=\sum_i a_i g_i$ with regular $a_i$.
Put $U=T_z\{g_1=g_2=g_3=0\}$, so $\dim U=6$.
For any $u,v,w\in U$, the product rule gives
\[
                    (D_uD_vD_w f)(z)=0.
\]
Indeed every summand contains either $g_i(z)$, a first derivative
$D_u g_i(z)$ with $u\in U$, or a second or higher derivative of $g_i$.
All of these vanish. The left side is the constant alternating
polar three-form of the cubic evaluated on $u,v,w$. Hence this form
vanishes on $\bigwedge^3 U$.

Choose a complement to $U$ in $V$. On $V^*$ give its three-dimensional
subspace $U^\perp$ weight $2$ and a complementary six-space weight $-1$.
These weights sum to zero. Every monomial of the polar form contains
at least one factor from $U^\perp$, so it has nonnegative weight.
The resulting nontrivial one-parameter subgroup has a finite limit
on that form. The Hilbert--Mumford stable-orbit criterion therefore
excludes stability. Equivalently, if the limit has only weight-zero
terms it has a positive-dimensional stabilizer, and if there are
positive-weight terms the original orbit is not closed with finite
stabilizer. This includes the case of zero polar form.
\end{proof}

\begin{theorem}[Full degree-one image under polar stability]
\label{thm:nonordinary-ppav-polar-stable-conditional-candidate}
Let $(A,\Theta)$ be a nonproduct nonordinary principally polarized
abelian surface over a perfect field $k$ of characteristic two.
Use its typed theta tensor and cubic $q=q_\alpha$ from
Lemma~\ref{lem:root-nonordinary-Coble-identification-candidate}.
Over $\overline k$ put
\[
 f(x)=q_{\alpha^{[1/2]}}(x),\qquad
 b_i=\left(\alpha_i\sum_{j\ne i}\alpha_j^3\right)^{1/2}.
\]
Suppose, as an additional hypothesis, that the alternating polar
three-form of $f$ is stable. In the affine-line theta basis its
four Cartan coefficients are precisely $b_i$.
Then the degree-one relative unit local system has rank eight and
both its geometric and arithmetic images equal $\GL_8(\mathbf F_2)$.
\end{theorem}
\begin{proof}
The rank is already established. For completeness we verify the
remaining exceptional-incidence implication. Write the nonzero
Lang equation as
\[
                 M(s)x^{[2]}=x
       \quad\Longleftrightarrow\quad M(x^{[2]})s=x.
\]
At a rank-six point choose an invertible principal alternating
six-by-six block. Solving for the first six coordinates with each
of the other three coordinates free gives three rational vectors
$k_i(w)$, independent throughout this block chart. On the rank-six
locus they form the kernel of $M(w)$. Since the matrix is alternating,
compatibility of the displayed linear equation in $s$ is exactly
\[
                    x^t k_i(x^{[2]})=0,\qquad i=1,2,3.
\]
The three functions have independent differentials. Their zero
scheme consequently has dimension six. Also the rank-six condition
implies $q(x^{[2]})=f(x)^2=0$.
If the rank-six compatible locus had a six-dimensional component,
its closure in this block chart would be a component of those three
smooth equations contained in $f=0$. The preceding lemma would make
the polar form of $f$ nonstable, contrary to the hypothesis. Thus
its dimension is at most five. Its affine solution spaces in $s$
have dimension three, so this entire exceptional incidence has
dimension at most eight.

Corollary~\ref{cor:root-nonordinary-Lang-rank-four-excluded-candidate}
already bounds the rank-four part by eight, excludes smaller nonzero
ranks, and identifies the rank-eight part as one irreducible
nine-dimensional affine-line torsor over an open of
$\sum x_i^3=0$. The projection to the section-parameter space is
\'etale, since the derivative in $x$ of $M(s)x^{[2]}-x$ is $-I$.
On the maximal-rank open the nonzero Lang scheme is finite \etale;
every component dominates that open and has dimension nine.
The bounds just proved therefore force it to have exactly one
component. Geometric monodromy is transitive on the 255 nonzero
vectors, in particular irreducible.

By Corollary~\ref{cor:nonordinary-ppav-uniform-root-inertia-candidate},
this group contains all transvections
\[
             1+u\otimes\varphi,\qquad u\in\ker\varphi,
\]
for some nonzero covector $\varphi$. The orbit of $\varphi$ spans
the whole dual space, since its annihilator is an invariant subspace
and the representation is irreducible. Choose an orbit basis
$\varphi_1,\ldots,\varphi_8$ and its dual basis $e_1,\ldots,e_8$.
Conjugating the displayed subgroup gives
$1+e_j\otimes\varphi_i$ for every $j\ne i$.
These elementary matrices generate $\GL_8(\mathbf F_2)$.
The arithmetic group contains the geometric group and is a subgroup
of the same general linear group, so equality holds for it as well.
\end{proof}

\begin{remark}[The new hypothesis does not follow from tensor stability]
\label{rem:nonordinary-polar-stability-gap-candidate}
This theorem is conditional. The extra hypothesis is verified for every
nonordinary principal surface in
Theorem~\ref{thm:every-nonordinary-ppav-degree-one-full-candidate} below;
the present argument alone does not verify it on the actual Cartier locus. In particular, stability of the original theta tensor does
not suffice by a formal invariant-theory implication.
In the four-dimensional Cartan coefficient space, consider
\[
                    \alpha_1^3+\alpha_2^3+\alpha_3^3=0.
\]
This is an irreducible cubic hypersurface: its plane base is a smooth
Fermat cubic in characteristic two. It is not contained in the finite
union of reflection hyperplanes. By
\citet[Proposition~2.10]{RainsSam2018Trivectors}, its general point
therefore defines a stable original tensor, with all $\alpha_i\ne0$.
But its polar coefficient $b_0$ is zero, so the polar tensor lies on
a coordinate reflection hyperplane and is nonstable by
\citet[Proposition~2.9]{RainsSam2018Trivectors}.
This proves the failure of the proposed implication on the entire
stable Cartan family. It does not identify that hypersurface with
any locus of actual Cartier tensors. The subsequent normal-form
calculations show that the actual nonordinary Cartier locus satisfies
polar stability, and thereby prove the unconditional degree-one
fixed-surface theorem without the false general implication.
\end{remark}
